\chapter{GIL, multiprocessing e async/await}

\section{GIL na pratica}

O GIL impede que multiplas threads executem bytecode Python puro simultaneamente no mesmo processo. Isso nao significa que threads sao inuteis. Significa que voce precisa entender o tipo de gargalo.

\begin{longtable}{p{4cm}p{8cm}}
\toprule
\textbf{Cenario} & \textbf{Ferramenta provavel} \\
\midrule
I/O-bound & threads ou async \\
CPU-bound & multiprocessing \\
Muitos arquivos pequenos & ThreadPoolExecutor pode ajudar \\
Transformacao numerica pesada & ProcessPoolExecutor pode ajudar \\
API web com muitas esperas & async/await pode ajudar \\
\bottomrule
\end{longtable}

\section{Multiprocessing}

Multiprocessing cria processos separados. Cada processo tem seu proprio interpretador Python e seu proprio GIL. Por isso ele pode acelerar tarefas CPU-bound.

Exemplo didatico:

\begin{lstlisting}[style=devbutterpython]
from concurrent.futures import ProcessPoolExecutor


def cpu_heavy_transform(value: float | None) -> float | None:
    if value is None:
        return None

    result = value
    for _ in range(100_000):
        result = (result * 1.000001) / 1.000001

    return result


def process_values(values: list[float | None]) -> list[float | None]:
    with ProcessPoolExecutor() as executor:
        return list(executor.map(cpu_heavy_transform, values))
\end{lstlisting}

\section{Custo do multiprocessing}

Multiprocessing nao e gratis:

\begin{itemize}
  \item cada processo consome mais memoria;
  \item dados precisam ser serializados entre processos;
  \item criar processos tem overhead;
  \item para tarefas pequenas, pode ficar mais lento.
\end{itemize}

Use multiprocessing quando o trabalho por item for pesado o suficiente.

\section{ThreadPoolExecutor para GPX}

Arquivos GPX sao bons candidatos para concorrencia porque geralmente existem varios arquivos menores.

\begin{lstlisting}[style=devbutterpython]
from concurrent.futures import ThreadPoolExecutor
from pathlib import Path


def parse_gpx_file(file: Path) -> list[dict[str, str | None]]:
    # parseia um arquivo GPX e retorna pontos
    return []


def parse_many_gpx(files: list[Path]) -> list[dict[str, str | None]]:
    rows: list[dict[str, str | None]] = []

    with ThreadPoolExecutor(max_workers=8) as executor:
        for file_rows in executor.map(parse_gpx_file, files):
            rows.extend(file_rows)

    return rows
\end{lstlisting}

\section{Async/await}

\code{async/await} e uma forma de concorrencia cooperativa. Em vez de criar varias threads, um event loop alterna entre tarefas quando elas fazem \code{await} em operacoes I/O.

Exemplo minimo:

\begin{lstlisting}[style=devbutterpython]
import asyncio


async def producer(queue: asyncio.Queue[int]) -> None:
    for item in range(10):
        await queue.put(item)

    await queue.put(-1)


async def consumer(queue: asyncio.Queue[int]) -> None:
    while True:
        item = await queue.get()

        if item == -1:
            break

        print(item)


async def main() -> None:
    queue: asyncio.Queue[int] = asyncio.Queue()

    await asyncio.gather(
        producer(queue),
        consumer(queue),
    )


asyncio.run(main())
\end{lstlisting}

\section{Onde async entra no projeto?}

Nao comece com async no parser principal. Primeiro faca a versao sincrona eficiente. Async pode entrar depois em:

\begin{itemize}
  \item API com FastAPI;
  \item consulta de status de parse;
  \item jobs em background;
  \item writer async com \code{asyncpg};
  \item dashboard consultando dados sem bloquear.
\end{itemize}

\section{Ordem recomendada}

\begin{enumerate}
  \item parser sincrono;
  \item logs e profiling;
  \item threads com queue;
  \item ThreadPool para GPX;
  \item ProcessPool para CPU-bound;
  \item async para API ou I/O concorrente futuro.
\end{enumerate}
